\documentclass[11pt,a4paper]{article}

\usepackage[margin=2.2cm]{geometry}
\usepackage{amsmath,amssymb,mathtools}
\usepackage{booktabs}
\usepackage{enumitem}
\usepackage{xcolor}
\usepackage{hyperref}
\usepackage{fancyhdr}
\usepackage{listings}
\usepackage{microtype}

\definecolor{ADGBlue}{RGB}{19,67,115}
\definecolor{ADGDark}{RGB}{22,29,38}
\definecolor{CodeBackground}{RGB}{245,247,249}
\definecolor{CodeGreen}{RGB}{20,110,60}

\hypersetup{
    colorlinks=true,
    hypertexnames=false,
    linkcolor=ADGBlue,
    urlcolor=ADGBlue,
    pdftitle={DisputeGuard Equations and Algorithms},
    pdfauthor={ADGSTUDIOS}
}

\pagestyle{fancy}
\fancyhf{}
\lhead{\textbf{ADGSTUDIOS}}
\rhead{DisputeGuard Equations and Algorithms}
\cfoot{\thepage}
\setlength{\headheight}{14pt}

\lstdefinestyle{pseudocode}{
    basicstyle=\ttfamily\small,
    backgroundcolor=\color{CodeBackground},
    frame=single,
    breaklines=true,
    breakatwhitespace=false,
    showstringspaces=false,
    numbers=left,
    numberstyle=\tiny,
    xleftmargin=1.5em,
    framexleftmargin=1em,
    commentstyle=\color{CodeGreen}
}

\begin{document}

\begin{titlepage}
    \centering
    \vspace*{2cm}
    {\Huge\bfseries ADGSTUDIOS\par}
    \vspace{0.5cm}
    {\LARGE\color{ADGBlue}\bfseries DisputeGuard\par}
    \vspace{1.5cm}
    \rule{0.85\textwidth}{0.8pt}
    \vspace{1cm}
    {\Large\bfseries Equations and Algorithms\par}
    \vspace{0.5cm}
    {\large Plain-language guide and implementation pseudocode\par}
    \vfill
    {\large Companion to the DisputeGuard Software Engineering Specification\par}
    \vspace{0.5cm}
    {\large Version 1.0 \quad|\quad October 2026\par}
    \vfill
    \textit{
        A practical reference for the calculations, safety rules and
        workflows used in dispute management.
    }
\end{titlepage}

\tableofcontents
\newpage

\section{How to Use This Guide}

This is a standalone companion to the DisputeGuard engineering
specification. It explains the main equations in everyday language and
gives pseudocode for important system workflows.

The pseudocode describes required behaviour, not a particular
programming language or database library. The main specification remains
the source of project requirements. Where a business rule is not yet
specified, this guide names the rule as a policy input instead of making
up a threshold or decision.

\subsection{Notation at a Glance}

\begin{center}
\begin{tabular}{@{}ll@{}}
\toprule
Symbol & Meaning \\
\midrule
\(u\) & A user \\
\(t\) & A financial transaction \\
\(d\) & A dispute about a transaction \\
\(s\) & A dispute state \\
\(e\) & An event or requested action \\
\(R(d)\) & A deterministic risk indicator for dispute \(d\) \\
\(\bot\) & An invalid or undefined result \\
\end{tabular}
\end{center}

\section{Money and Exact Amounts}

For a currency that has \(q\) minor-unit decimal places, let
\(e_c=10^q\). An amount \(a\), parsed as an exact decimal, is represented
in integer minor units as:

\[
a_{\mathrm{minor}} = a \times e_c.
\]

For a two-decimal currency:

\[
\text{R}1234.56 \longrightarrow 123456\text{ cents}.
\]

\textbf{In plain language:} Store cents as whole numbers instead of
binary floating-point numbers. This prevents rounding surprises when
adding or comparing money. Use the decimal-place rule for the actual
currency; do not assume every currency has two decimal places. Only add
amounts after confirming they use the same currency.

\section{Ownership and Permission}

Let \(Owner(u,t)\) mean that user \(u\) owns the account associated with
transaction \(t\). A customer may access their own transaction; an
analyst or administrator may access it only when their role grants that
permission:

\[
CanView(u,t)
\iff
Owner(u,t)
\lor
Role(u)\in\{\text{ANALYST},\text{ADMIN}\}.
\]

For customer-created disputes, ownership is mandatory:

\[
CanCreateDispute(u,t)
\Rightarrow
Owner(u,t).
\]

\textbf{In plain language:} A transaction ID supplied by a browser is
not proof that the customer owns that transaction. The server checks
ownership using the authenticated user's identity on every request.
Hiding a button in the user interface is not an access-control check.

\section{Dispute Lifecycle}

Let \(S\) be the configured set of dispute states and \(E\) the set of
possible events. The transition function is:

\[
T:S\times E\rightarrow S\cup\{\bot\}.
\]

Some transitions specified for this system are:

\begin{align*}
T(\text{SUBMITTED},\text{BEGIN\_REVIEW})
    &= \text{UNDER\_REVIEW},\\
T(\text{UNDER\_REVIEW},\text{REQUEST\_EVIDENCE})
    &= \text{EVIDENCE\_REQUIRED},\\
T(\text{EVIDENCE\_REQUIRED},\text{EVIDENCE\_SUBMITTED})
    &= \text{EVIDENCE\_RECEIVED},\\
T(\text{INVESTIGATING},\text{APPROVE})
    &= \text{APPROVED},\\
T(\text{INVESTIGATING},\text{REJECT})
    &= \text{REJECTED}.
\end{align*}

An action not present in the authorised transition table is invalid.
For example:

\[
T(\text{SUBMITTED},\text{ISSUE\_REFUND})=\bot.
\]

\textbf{In plain language:} A case can only move through steps that the
workflow allows. Every attempted change is checked by the server; the
user interface cannot skip a review step. The full transition table and
role permissions must be configured by the application.

\section{One Active Dispute per Transaction}

Let \(S_A\) be the states considered active by the application. For
every transaction \(t\), the number of active disputes is limited to
one:

\[
\forall t,\quad
\left|
\left\{
d : d.transaction=t\land d.status\in S_A
\right\}
\right|
\leq 1.
\]

\textbf{In plain language:} Two requests arriving at nearly the same
time must not create two open cases for the same transaction. The
database must enforce this rule; an application-level check alone can
lose a race.

\section{Safe, Repeatable Requests}

For an idempotency key \(K\), authenticated user \(u\), canonical
request hash \(h\), and original result \(r\), the server records:

\[
(u,K,h)\mapsto r.
\]

On a retry:

\[
\begin{cases}
\text{return }r,
    &\text{if }u,K,h\text{ match a completed request},\\
\text{reject},
    &\text{if the key is reused with a different request},\\
\text{process once},
    &\text{if no matching request exists}.
\end{cases}
\]

\textbf{In plain language:} If a customer clicks submit twice or the
network retries a request, the customer gets the original result rather
than a duplicate dispute. Reusing the same key for different request
contents must not silently return or create the wrong result.

\section{Concurrent Updates}

Each dispute has an integer version \(v\). A change succeeds only when
the version the caller saw still matches the database version:

\[
v_{\mathrm{request}}=v_{\mathrm{database}}.
\]

On success:

\[
v_{\mathrm{new}}=v_{\mathrm{old}}+1.
\]

\textbf{In plain language:} If two analysts edit a case at once, the
second save must notice that the case changed since it was opened. It
should ask the user to reload or resolve the conflict, not overwrite a
newer decision.

\section{Deterministic Risk Indicators}

Let \(A,V,M,H,G\in[0,1]\) be available, normalised signals for amount,
velocity, merchant, historical behaviour and geography. Let the
configured weights be non-negative and sum to one:

\[
w_A+w_V+w_M+w_H+w_G=1,
\qquad
w_A,w_V,w_M,w_H,w_G\geq0.
\]

The weighted score is:

\[
R(d)=w_AA(d)+w_VV(d)+w_MM(d)+w_HH(d)+w_GG(d).
\]

When every input is in \([0,1]\), this implies:

\[
0\leq R(d)\leq1.
\]

\textbf{In plain language:} The score combines several transparent
signals using configured weights. It helps an analyst decide what to
look at; it does not prove fraud and must not approve or reject a
dispute. A missing signal is unknown, not zero. Do not publish a
combined score when required inputs are missing unless an explicit,
documented missing-data policy says how to handle them.

\subsection{Amount Normalisation}

For a non-negative amount \(x\) and a positive configured scale
\(x_{\max}\):

\[
A(x)=
\min\left(1,\frac{\ln(1+x)}{\ln(1+x_{\max})}\right).
\]

\textbf{In plain language:} This compresses very large values so one
unusually large purchase cannot automatically dominate the other risk
signals. The scale must be configured for the relevant currency and
population.

\subsection{Historical Amount Anomaly}

For \(n\geq2\) prior amounts \(x_1,\ldots,x_n\), the sample mean and
standard deviation are:

\[
\bar{x}=\frac{1}{n}\sum_{i=1}^{n}x_i,
\qquad
s=\sqrt{\frac{1}{n-1}\sum_{i=1}^{n}(x_i-\bar{x})^2}.
\]

When \(s>0\), the standardised deviation for a new amount \(x\) is:

\[
z=\frac{x-\bar{x}}{s}.
\]

An absolute \(z\)-score above \(3\) may be labelled unusual for
demonstration purposes; it is not evidence that fraud occurred. If
there is insufficient history or \(s=0\), the result is unavailable
rather than an invented score.

\textbf{In plain language:} Compare a new amount with the customer's
past amounts. A value far from the usual range is a reason to take a
closer look, not a verdict.

\subsection{Robust Amount Anomaly}

To reduce the influence of extreme past values, define the median \(m\),
the median absolute deviation \(MAD\), and, when \(MAD>0\):

\[
m=\operatorname{median}(x_1,\ldots,x_n),
\qquad
MAD=\operatorname{median}(|x_i-m|),
\]

\[
z_i^*=\frac{0.6745(x_i-m)}{MAD}.
\]

\textbf{In plain language:} The median-based calculation is less
affected by a few unusually large historical purchases. If \(MAD=0\),
the calculation is undefined; do not divide by zero.

\subsection{Transaction Velocity}

For a window of length \(\Delta t>0\), let
\(N_u(t-\Delta t,t]\) count a customer's transactions in that window:

\[
V_u(t)=\frac{N_u(t-\Delta t,t]}{\Delta t}.
\]

When historical velocity has standard deviation
\(\sigma_{V,u}>0\), compare current velocity with history:

\[
z_V=\frac{V_u(t)-\mu_{V,u}}{\sigma_{V,u}}.
\]

\textbf{In plain language:} Count transactions in a defined time
window, then compare that rate with the customer's normal rate. The
window and any threshold must be explicit configuration, not an
unstated assumption.

\subsection{Geographical Signal}

For two latitude/longitude points in radians, let
\(\Delta\phi=\phi_2-\phi_1\) and
\(\Delta\lambda=\lambda_2-\lambda_1\). The Haversine calculation is:

\begin{align*}
a &=
\sin^2\left(\frac{\Delta\phi}{2}\right)
+\cos(\phi_1)\cos(\phi_2)
\sin^2\left(\frac{\Delta\lambda}{2}\right),\\
c &=2\operatorname{atan2}\left(\sqrt{a},\sqrt{1-a}\right),\\
d &=R_Ec,
\qquad R_E\approx6371\text{ km},\\
v_{\mathrm{required}}&=\frac{d}{\Delta t},
\qquad \Delta t>0.
\end{align*}

\textbf{In plain language:} Estimate how far apart two recorded
locations are, then how quickly someone would need to travel between
them. Missing locations, non-positive elapsed time, online purchases,
VPN use and delayed settlement make this signal uncertain. It is never
proof of fraud by itself.

\section{Evidence Upload Workflow}

\textbf{In plain language:} Treat a file as untrusted until the server
has checked who uploaded it, how large it is, what it contains and
whether scanning is required and available. A filename or browser
MIME-type claim is not proof of file content.

\begin{lstlisting}[style=pseudocode,caption={Validating and recording evidence}]
function receiveEvidence(user, disputeId, upload):
    authenticate(user)
    dispute = loadDispute(disputeId)
    require canAccessAndSubmitEvidence(user, dispute)

    require upload.size <= configuredSizeLimit
    require upload.declaredType in configuredAllowedTypes
    require extensionAllowed(upload.name)
    require detectedContentType(upload.bytes) in configuredAllowedTypes

    storageName = generateRandomServerSideName()
    storeOutsideExecutableWebPaths(storageName, upload.bytes)

    if malwareScanningIsRequired:
        require scannerIsAvailable()
        require scan(storageName) == CLEAN

    begin database transaction
        saveEvidenceMetadata(dispute, storageName, upload)
        appendEvidenceAuditEvent(user, dispute)
    commit

    return evidenceReceipt()
\end{lstlisting}

If storage or scanning succeeds but metadata persistence fails, the
implementation must clean up or safely reconcile the unreferenced
object. Scanner unavailability must not be reported as a clean scan.

\section{Dispute Creation Algorithm}

\textbf{In plain language:} Verify the customer's identity and
ownership, apply the configured eligibility rules, and create the case
and its audit records as one database operation. The database, not a
pre-check in the interface, protects against simultaneous duplicate
submissions.

\begin{lstlisting}[style=pseudocode,caption={Creating a dispute safely}]
function createDispute(authenticatedUser, transactionId, request, key):
    require authenticatedUser is authenticated
    canonicalRequest = canonicalize(request, transactionId)
    requestHash = hash(canonicalRequest)

    begin database transaction

    replay = findIdempotentResult(authenticatedUser.id, key)
    if replay exists:
        if replay.requestHash != requestHash:
            rollback
            reject key reuse with different request
        commit
        return replay.originalResult

    transaction = loadTransactionForUpdate(transactionId)
    require transaction exists
    require transaction.customerId == authenticatedUser.id
    require eligibilityPolicyAllows(transaction, request.reason)
    require databaseAllowsNoOtherActiveDispute(transactionId)

    dispute = insertDispute(transaction, authenticatedUser, request)
    appendDisputeEvent(dispute, "SUBMITTED", authenticatedUser)
    appendAuditEvent("DISPUTE_CREATED", authenticatedUser, dispute)
    saveIdempotentResult(authenticatedUser.id, key, requestHash, dispute)

    commit
    return dispute

on any failure:
    rollback all writes
    return an explicit error
\end{lstlisting}

The idempotency record and dispute must be committed atomically. The
transaction lock and database uniqueness constraint are both needed:
the lock serialises the competing operation, while the constraint
remains the final guard against races.

\section{Authorised State Transition Algorithm}

\textbf{In plain language:} Load the latest case, confirm the actor is
allowed to perform the requested action, check that the action is valid
from the current state, and save the new state together with its audit
record. Reject stale or invalid changes.

\begin{lstlisting}[style=pseudocode,caption={Changing a dispute state}]
function transitionDispute(actor, disputeId, event, expectedVersion):
    require actor is authenticated
    begin database transaction

    dispute = loadDisputeForUpdate(disputeId)
    require actorMayPerform(actor, event, dispute)
    require dispute.version == expectedVersion

    nextState = authorisedTransition(dispute.state, event)
    require nextState exists

    oldState = dispute.state
    update dispute:
        state = nextState
        version = version + 1
    appendDisputeEvent(dispute, oldState, nextState, actor, event)
    appendAuditEvent("DISPUTE_TRANSITION", actor, dispute)

    commit
    return updated dispute

on any failed check:
    rollback
    return an explicit error
\end{lstlisting}

\section{Risk Calculation Algorithm}

\textbf{In plain language:} Calculate each signal with its own defined
inputs and safe mathematical guards. If a necessary input is missing or
undefined, show that fact rather than quietly treating it as a
low-risk value.

\begin{lstlisting}[style=pseudocode,caption={Calculating an explainable risk indicator}]
function calculateRisk(dispute):
    signals = {
        amount: calculateAmountSignal(dispute),
        velocity: calculateVelocitySignal(dispute),
        merchant: calculateMerchantSignal(dispute),
        history: calculateHistorySignal(dispute),
        geography: calculateGeographySignal(dispute)
    }

    if any required signal is unavailable:
        return {
            score: unavailable,
            signals: signals,
            explanation: "One or more required inputs are unavailable"
        }

    require every signal is in [0, 1]
    require configured weights are nonnegative
    require sum(configured weights) == 1

    score = sum(weight[name] * signals[name] for each signal)
    return {score: score, signals: signals}
\end{lstlisting}

Risk bands and thresholds are configuration for analyst triage. The
result is not a fraud finding or an automated financial decision.

\section{Bounded AI Investigation Algorithm}

\textbf{In plain language:} An AI helper may gather a limited amount of
approved information and suggest a summary. It cannot grant itself new
tools, run arbitrary database queries, or make the final decision. A
human with the right permissions remains responsible for case outcomes.

\begin{lstlisting}[style=pseudocode,caption={Running a bounded investigation}]
function investigate(caseId, configuredBudget):
    context = loadOnlyAuthorisedCaseContext(caseId)
    state = initialiseInvestigation(context, configuredBudget)

    while not state.complete:
        require state.iterations < configuredBudget.maxIterations
        require state.toolCalls < configuredBudget.maxToolCalls
        require state.tokens < configuredBudget.maxTokens
        require elapsedTime(state) < configuredBudget.maxDuration

        proposal = requestNextStructuredStep(state)
        require validateSchema(proposal)
        require proposal.tool in applicationAllowlist

        result = executeReadOnlyToolForThisCase(proposal)
        state = recordStructuredObservation(state, proposal, result)

    recommendation = validateSchema(requestFinalSummary(state))
    persistSummaryAndToolHistory(recommendation, state)
    return recommendation for authorised human review

on budget exhausted, invalid output, or model failure:
    record analysis as incomplete or failed
    do not change dispute state or financial balances
\end{lstlisting}

Customer descriptions and uploaded evidence are data, not instructions
that can grant authority. The application enforces tool permissions and
execution bounds independently of model prompts.

\section{Safe Background Retries}

\textbf{In plain language:} A worker can receive the same job more than
once after a timeout or crash. Before retrying, the application must
know that repeating the operation cannot duplicate an external effect.

\begin{lstlisting}[style=pseudocode,caption={Handling retryable background work}]
function handleJob(job):
    require job.operation is explicitly retryable
    require job has a stable idempotency identifier

    if durableResultExists(job.idempotencyIdentifier):
        acknowledge job
        return stored result

    try:
        result = performOperationWithIdempotency(job)
        saveDurableResult(job.idempotencyIdentifier, result)
        acknowledge job
    catch retryableFailure:
        retry only while configured retry policy permits
        otherwise record failure for operator review
    catch nonRetryableFailure:
        record failure for operator review
        do not retry automatically
\end{lstlisting}

Exactly-once effects across independent services cannot be assumed from
a queue alone. Durable idempotency records, transactional writes and
explicit failure handling are needed.

\section{Key Safety Invariants}

\begin{center}
\begin{tabular}{@{}p{0.32\textwidth}p{0.58\textwidth}@{}}
\toprule
Invariant & What it means in practice \\
\midrule
Money is exact & Use integer minor units or fixed-precision decimal values. \\
Ownership is server-checked & A caller cannot gain access by changing an identifier. \\
At most one active dispute & Database constraints protect simultaneous submissions. \\
Transitions are authorised & Only configured actor/event/state combinations can proceed. \\
Retries are safe & A repeated request returns its original result, not a duplicate effect. \\
Risk is advisory & A score or anomaly is not proof and cannot decide a case. \\
AI is bounded & Allowlisted tools and resource limits are enforced by application code. \\
Audit is durable & Case changes and their audit records are saved together. \\
\bottomrule
\end{tabular}
\end{center}

\end{document}
